% !TeX root = ../main.tex
%------------------------------------------------------------------------------
%  Appendix E -- Stereographic Net Constructions
%  Procedures only. The concepts and the underground application are in
%  Chapter~\ref{ch:structural-analysis}.
%------------------------------------------------------------------------------
\chapter{Stereographic Net Constructions}
\label{app:stereonet}

Each construction below is set out as numbered steps that can be followed on a
printed net with tracing paper, so that the output of stereonet software can
be checked by hand. The reasoning behind them is in
Chapter~\ref{ch:structural-analysis}. All constructions assume lower
hemisphere projection and the dip/dip-direction convention.

\section{Setting up: nets, tracing paper and conventions}
\label{sec:net-setup}

Two nets are used and they are not interchangeable. The \emph{equal-angle}
(Wulff) net preserves angular relationships and is used for all geometric
construction --- intersections, angles, rotations. The \emph{equal-area}
(Schmidt) net preserves area and is used for pole density contouring and any
statistical work. Using an equal-angle net for contouring concentrates
apparent density towards the centre of the plot and will invent a set that
does not exist.

\begin{enumerate}
  \item Fix the net to a board and place a sheet of tracing paper over it,
        pinned through the centre so that it rotates freely.
  \item Trace the primitive circle and mark north on the tracing paper. This
        index mark is the reference for every subsequent rotation, and all
        readings are taken with it returned to the top of the net.
  \item Record the convention on the sheet: lower hemisphere, dip/dip
        direction, and the magnetic declination correction applied to the
        field data if any.
\end{enumerate}

Convention errors are the commonest source of gross mistakes. Data recorded as
strike and dip must be converted before plotting --- dip direction is strike
plus $90^\circ$ under the right-hand rule --- and mixing the two produces
planes rotated by a quarter turn.

\section{Plotting planes, poles and lines}
\label{sec:plotting-planes-poles-and-lines}

\subsection{Plotting a plane as a great circle}
\label{sec:construct-great-circle}

Given a plane with dip $\psi$ and dip direction $\alpha$:

\begin{enumerate}
  \item Rotate the tracing paper until the index mark reads $\alpha$
        clockwise from north on the net.
  \item Count $\psi$ degrees inward from the primitive along the
        east--west diameter of the net.
  \item Trace the great circle passing through that point.
  \item Return the index mark to north.
\end{enumerate}

\noindent\textbf{Check.}\enspace A vertical plane plots as a straight diameter;
a horizontal plane plots as the primitive circle itself. If the traced great
circle bulges the wrong way, the index was rotated in the wrong sense.

\subsection{Plotting the pole to a plane}
\label{sec:construct-pole}

The pole is the projection of the normal to the plane, and it is a single
point --- which is why poles rather than great circles are used when many
planes must be plotted together.

\begin{enumerate}
  \item With the index at the dip direction $\alpha$ as above, count
        $\psi$ inward from the primitive to locate the great circle.
  \item Continue counting a further $90^\circ$ along the same diameter,
        passing through the centre. Mark that point: it is the pole.
  \item Return the index to north.
\end{enumerate}

Equivalently, the pole may be plotted directly by treating it as a line of
plunge $90^\circ - \psi$ and trend $\alpha + 180^\circ$.

\noindent\textbf{Check.}\enspace The pole to a horizontal plane is at the
centre of the net; the pole to a vertical plane lies on the primitive.

\subsection{Plotting a line from plunge and trend}
\label{sec:construct-line}

Given a line of plunge $\beta$ and trend $\theta$ --- a lineation, a borehole
axis, or an intersection:

\begin{enumerate}
  \item Rotate the index to read $\theta$.
  \item Count $\beta$ degrees inward from the primitive along the
        east--west diameter.
  \item Mark the point and return the index to north.
\end{enumerate}

To recover plunge and trend from a plotted point, reverse the procedure:
rotate until the point lies on the east--west diameter, read the trend at the
index and the plunge as the angular distance from the primitive.

\section{Intersections, angles and rotation}
\label{sec:intersections-angles-and-rotation}

\subsection{Intersection of two planes}
\label{sec:construct-intersection}

This is the construction underlying every wedge analysis: the line of
intersection is the direction along which a wedge formed by the two planes
will slide.

\begin{enumerate}
  \item Plot both planes as great circles.
  \item Mark the point where the two great circles cross. This point is the
        line of intersection.
  \item Rotate the tracing until that point lies on the east--west diameter.
  \item Read the trend at the index mark and the plunge as the angular
        distance inward from the primitive.
\end{enumerate}

\noindent\textbf{Check.}\enspace The intersection of two planes must lie on
both great circles, and its plunge cannot exceed the dip of either plane. A
plunge steeper than both dips indicates a plotting error.

\subsection{Angle between two lines or two planes}
\label{sec:construct-angles}

\begin{enumerate}
  \item For two lines, plot both as points, then rotate the tracing until
        both lie on the same great circle.
  \item Count the angular separation along that great circle. This is the
        true angle between them.
  \item For two planes, plot their poles and measure the angle between the
        poles by the same method. The dihedral angle between the planes is
        the supplement of the angle between their poles.
\end{enumerate}

Angles must be counted along a great circle, never as a straight-line distance
on the projection, which is distorted.

\subsection{Rotation about an axis}
\label{sec:construct-rotation}

Rotation is used to correct data for a plunging fold axis, to reorient data
from an inclined borehole, and to test alternative excavation alignments.

\begin{enumerate}
  \item Plot the rotation axis and rotate the tracing until it lies at the
        centre of the net, or on the primitive for a horizontal axis.
  \item Move every plotted point along its small circle, through the required
        angle, in the required sense.
  \item Restore the tracing to its original position and read the rotated
        orientations.
\end{enumerate}

For a horizontal axis the small circles are those of the net directly. For an
inclined axis, the data must first be rotated to bring the axis to a
convenient position, the required rotation applied, and the first rotation
reversed.

\section{Statistical treatment of orientation data}
\label{sec:statistical-treatment-of-orientation-data}

\paragraph{Pole density contouring.}
Contouring identifies discontinuity sets within a scatter of poles. It must be
done on an equal-area net.

\begin{enumerate}
  \item Plot all poles on the equal-area net.
  \item Overlay a counting grid. Using a counting circle of one per cent of
        the net area, count the poles falling within the circle centred at
        each grid node; at the primitive, use the paired counting circle so
        that poles near the edge are counted once.
  \item Express each count as a percentage of the total number of poles.
  \item Contour the resulting values at intervals chosen so that three to
        five contours span the range.
\end{enumerate}

\noindent\textbf{Minimum data.}\enspace Contouring fewer than about a hundred
poles produces contours that describe the sample rather than the rock mass. A
plot of thirty poles should be presented as poles, not as contours.

\paragraph{Terzaghi bias correction.}
A scanline or borehole under-samples discontinuities lying close to parallel
with it. The correction weights each observation by

\[
w = \frac{1}{\sin\delta}
\]

where $\delta$ is the acute angle between the discontinuity plane and the
sampling line.

\begin{enumerate}
  \item For each observation, compute $\delta$ from the plotted pole and the
        plotted sampling line.
  \item Apply the weight $w$, counting the observation $w$ times in the
        density calculation.
  \item Apply a cut-off: where $\delta$ falls below a chosen minimum,
        conventionally $15^\circ$ to $20^\circ$, cap $w$ at the value
        corresponding to the cut-off.
\end{enumerate}

The cut-off is necessary because $w$ diverges as $\delta \to 0$, and without
it a single observation of a near-parallel set is multiplied into an apparent
cluster. Sets close to parallel with the sampling line remain
under-represented after correction, and the only remedy is to sample in more
than one direction.

\section{Block and wedge construction}
\label{sec:block-and-wedge-construction}

\paragraph{Joint pyramid and removability.}

\begin{enumerate}
  \item Plot the great circles of the discontinuity sets bounding the block.
        These divide the net into regions.
  \item Select the region corresponding to the block under consideration.
        This region is the joint pyramid: the set of directions in which the
        block is free to move as determined by its faces.
  \item Plot the great circle of the excavation surface, and identify the
        region representing directions into the rock. This is the excavation
        pyramid.
  \item Apply the emptiness test. If the joint pyramid and the excavation
        pyramid have no region in common, the block is removable. If they
        overlap, the block is constrained by the surrounding rock and cannot
        move regardless of its weight.
\end{enumerate}

\noindent\textbf{Check.}\enspace Removability depends only on orientation. A
block that fails the test cannot move, and no calculation of its weight is
required.

\paragraph{Wedge geometry from the projection.}
For a removable tetrahedral wedge formed by three planes in the crown:

\begin{enumerate}
  \item Determine the three lines of intersection of the bounding planes,
        taken in pairs.
  \item Identify the sliding direction: the line of intersection that
        daylights into the excavation for a two-plane slide, or the dip
        vector of the single sliding plane for a one-plane slide. If no
        intersection daylights and the block is in the crown, the mode is
        free fall.
  \item Construct the wedge in section, using the excavation span and the
        plane orientations, to obtain the apex height.
  \item Compute the base area from the excavation geometry and the trace
        lengths, the volume as one third of base area times apex height, and
        the weight as volume times unit weight.
\end{enumerate}

The stability calculation that follows is in
Appendix~\ref{app:calculations}.

\section{Blank nets}
\label{sec:blank-nets}

Printable equal-angle and equal-area nets follow, at a diameter suitable for
hand construction with a protractor. Both are drawn at two-degree intervals.
